\section{Introduction}
A quantum instrument can be described once and then used repeatedly.
We ask how many classical description bits suffice to approximate
all interactions with that instrument, uniformly over a common
causal quantum tester. The description specifies a genuine quantum
instrument; it does not replace that instrument by a classical
physical device. The target is supplied to the encoder, rather than
learned from queries.

For an $m$-outcome instrument from $M_d$ to $M_n$, write
$s=d^2(mn^2-1)$. Fix $d,n,m$ and a blockwise Choi eigenvalue margin
$0<a<1/(mn)$. Theorem~\ref{thm:jointcode68} proves
\begin{equation*}
 \log_2\mathcal M_N(\delta;a)
 =\frac{s}{2}\log_2 N+s\log_2(1/\delta)+O_{d,n,m,a}(1),
 \qquad N\ge1,\quad 0<\delta\le1/32.
\end{equation*}
The metric is the unhalved final-state trace distance, maximized over
common causal testers with entangled references, arbitrary finite
quantum memory, feedback and bounded public stopping. The same
memoryless instrument is reused at every invocation. The remainder
is independent of both $N$ and $\delta$.

A complementary theorem includes a singular boundary. An
outcome-retaining preparation instrument sends $X$ to
$(\tr(X)\sigma_y)_y$, where $\sigma_y\ge0$ and
$\sum_y\tr\sigma_y=1$. For fixed rank bounds $r_y$, put
\[
                  v=\sum_y r_y(2n-r_y)-1.
\]
Theorem~\ref{thm:rankentropy68} gives
\[
 \log_2\mathcal P_N(\delta;\mathbf r)
 =\frac v2\log_2N+v\log_2(1/\delta)+O_{n,m,\mathbf r}(1)
\]
on the same joint range, including zero outcomes and rank-deficient
blocks, without any eigenvalue floor. The rational code of
Theorem~\ref{thm:factorcode68} preserves zero outcomes and never
increases any block rank. Its integer square roots encode ratios of
triangular-factor coordinates; the encoder need not store algebraic
numbers. The singleton case $v=0$ has zero payload length.

The joint converse rests on a binary product-testing bound that is
linear in the local separation down to arbitrarily small errors.
An affine packing then retains the term $s\log_2(1/\delta)$, which
would be hidden by a fixed-error remainder. On a positive Choi body,
the upper bound uses a two-point classical programme. For preparation
instruments, input erasure instead reduces the full adaptive metric
exactly to the trace distance between product states. Purification
and rank-preserving factor coordinates give a boundary-uniform upper
code. Volume, binary testing, programme arguments and purification
are classical ingredients; their roles are separated in the proofs.

The single-use intrinsic entropy and coordinate codec, and the
fixed-error adaptive law, are retained as complete preliminary
results. No action spectral gap is used in these description theorems.
The complete research edition separately retains the spectral width,
positive numerical streaming and structural realization theory, with
all their original hypotheses and proofs. Neither that larger theory
nor its historical archive is a prerequisite for the present article.

Dimensions, named outcomes, and the requested public parameters are
fixed in the leading payload laws. A rank bound is also public where
used; the finite pivot mask is charged in the explicit factor-code
bound. Leading zero padding is counted conceptually in a fixed-length
code even though the hexadecimal interchange format is canonical and
unpadded. Headers, validation workspace and decoded matrix storage
are additional resources when their sizes vary. Decoding a word can
verify legality, not closeness to an unspecified target. The classical
instrument codec uses its validated encoder image, not its whole
mixed-radix cube, as a legal codebook.
